% Auto-generated from meeting_2567-08-16_pre-launch-preparation.md (กบ.วช. format v2)
\documentclass{meetingminutes}
\meetingcommittee{คณะกรรมการบริหารหลักสูตร วิศวกรรมคอมพิวเตอร์และปัญญาประดิษฐ์ (บัณฑิตศึกษา)}
\meetingnumber{๔}{๒๕๖๗}
\meetingdate{วันที่ ๑๖ สิงหาคม พ.ศ. ๒๕๖๗}
\meetingplace{ห้องประชุมคณะเทคโนโลยีอุตสาหกรรม}
\meetingstart{๐๙.๐๐ น.}
\meetingend{๑๑.๒๐ น.}

\begin{document}
\maketitle

\psection{ผู้มาประชุม}
\begin{participants}
  \pmember{1}{รองศาสตราจารย์ ดร.อิสระ อินจันทร์}{ที่ปรึกษาหลักสูตร}{ประธาน}
  \pmember{2}{รองศาสตราจารย์ ดร.กันต์ อินทุวงศ์}{ที่ปรึกษาหลักสูตร}{รองประธาน}
  \pmember{3}{ผู้ช่วยศาสตราจารย์ ดร.พิศิษฐ์ นาคใจ}{อาจารย์ประจำหลักสูตร}{กรรมการ}
  \pmember{4}{ผู้ช่วยศาสตราจารย์ ดร.กาญจนา ดาวเด่น}{อาจารย์ประจำหลักสูตร (เลขานุการ)}{กรรมการ/เลขานุการ}
  \pmember{5}{ผู้ช่วยศาสตราจารย์ ดร.พัชรี มณีรัตน์}{อาจารย์ประจำหลักสูตร}{กรรมการ}
  \pmember{6}{ผู้ช่วยศาสตราจารย์ ดร.พิทักษ์ คล้ายชม}{รองคณบดีฝ่ายวิชาการ คณะเทคโนโลยีอุตสาหกรรม}{กรรมการ}
  \pmember{7}{ผู้ช่วยศาสตราจารย์ ดร.โสกณ วิริยะรัตนกุล}{อาจารย์ประจำหลักสูตร}{กรรมการ}
  \pmember{8}{ผู้ช่วยศาสตราจารย์ ดร.ธัชชัย อยู่ยิ่ง}{อาจารย์ประจำหลักสูตร}{กรรมการ}
  \pmember{9}{ผู้ช่วยศาสตราจารย์ ดร.อภิศักดิ์ พรหมฝ่าย}{อาจารย์ประจำหลักสูตร}{กรรมการ}
\end{participants}

\psection{ผู้ไม่มาประชุม}
\noindent ไม่มี\par

\startmeeting
\openingchair{รองศาสตราจารย์ ดร.อิสระ อินจันทร์}{กล่าวเปิดการประชุมและดำเนินการประชุมตามระเบียบวาระ ดังนี้}

\agenda{๑}{รับรองรายงานการประชุมครั้งที่ 3/2567}
ที่ประชุมรับรองรายงานการประชุมครั้งที่ 3/2567 โดยไม่มีการแก้ไข


\agenda{๒}{เตรียมความพร้อมด้านอาจารย์และสิ่งอำนวยความสะดวก}
\subagenda{๒.๑}{ด้านอาจารย์ประจำหลักสูตร}
ที่ประชุมทบทวนรายชื่ออาจารย์ประจำหลักสูตร 7 คน และอาจารย์ผู้สอนเพิ่มเติม 6 คน ยืนยันว่าทุกคนมีคุณสมบัติครบถ้วนตามเกณฑ์ กมอ. (ปร.ด./วศ.ด./วท.ด. + ตำแหน่ง ผศ./รศ.) และมีภาระงานสอนไม่เกินขีดจำกัดที่กำหนด

ที่ประชุมพิจารณาการมอบหมายรายวิชาให้อาจารย์แต่ละท่านในภาคเรียนที่ 2/2568 (ภาคแรกที่เปิดสอน)

\subagenda{๒.๒}{ด้านสิ่งอำนวยความสะดวกและทรัพยากร}
รายงานความพร้อมของ: - ห้อง Lab คอมพิวเตอร์ 7 ห้อง (ชั้น 2 อาคาร ICIT) — พร้อมใช้งาน - คอมพิวเตอร์ 277 เครื่อง ระบบ GPU workstation และ Cloud computing - ระบบ WiFi 795 จุด ครอบคลุมพื้นที่เรียน - ฐานข้อมูลออนไลน์ 8 แห่ง (TCI, Scopus, IEEE Xplore ฯลฯ)


\agenda{๓}{วางแผนการประชาสัมพันธ์และรับสมัครนักศึกษา}
ที่ประชุมวางแผนการเปิดรับนักศึกษาปีการศึกษา 2568 ภาคเรียนที่ 2: - ช่องทางประชาสัมพันธ์: Facebook Page หลักสูตร, เว็บไซต์ มรอ., สื่อสิ่งพิมพ์ - เกณฑ์รับนักศึกษา: สำเร็จการศึกษาระดับปริญญาตรีสาขาวิศวกรรมคอมพิวเตอร์/วิทยาการคอมพิวเตอร์/หรือสาขาที่เกี่ยวข้อง - กำหนดรับสมัคร: ตุลาคม–พฤศจิกายน 2568 - ขั้นตอน: สมัคร → สอบสัมภาษณ์ → ประกาศผล → รายงานตัว

\resolution{เห็นชอบแผนการรับสมัครและมอบหมายให้ ผศ.ดร.พิศิษฐ์ เป็นประธานคณะกรรมการรับสมัครนักศึกษา}

\agenda{๔}{กำหนดการจัดทำ มคอ.3 รายวิชาภาคเรียนที่ 2/2568}
ที่ประชุมกำหนดตารางการจัดทำ มคอ.3 สำหรับทุกวิชาที่จะเปิดสอนในภาคเรียนที่ 2/2568 โดยมอบหมายให้อาจารย์แต่ละท่านส่ง มคอ.3 ภายในวันที่ 31 ตุลาคม 2568


\adjournmeeting

\begin{sigblock}
\typist{ผู้ช่วยศาสตราจารย์ ดร.กาญจนา ดาวเด่น}
\reviewer{รองศาสตราจารย์ ดร.อิสระ อินจันทร์}{กรรมการ}
\chairsig{รองศาสตราจารย์ ดร.อิสระ อินจันทร์}{ประธานการประชุม}
\end{sigblock}

\end{document}
